\documentclass[a4paper,10pt]{article}

\usepackage[left=1.4cm,top=1.2cm,bottom=1.5cm,right=1.2cm]{geometry}%

\usepackage[latin1]{inputenc}
\usepackage{amsmath}
%\usepackage{babel}
\topmargin 0truecm
\textheight 23truecm
\textwidth 16.5truecm
\oddsidemargin 0truecm
\parskip 1truemm
%\parsep 1truemm
\parindent 0em
\usepackage{amssymb}

\def\d{\,{\rm d}}

\def\RR{{\mathbb R}}
\def\CC{{\mathbb C}}
\def\NN{{\mathbb N}}
\def\ZZ{{\mathbb Z}}
\def\EE{{\mathbb E}}
\newcounter{problema}
\newcommand{\prob}[1]{\vspace{0.33cm}\stepcounter{problema}
                 \noindent{\bf Problema \arabic{problema}:}{\it #1}}

\begin{document}
\renewcommand{\labelenumi}{\bf \alph{enumi})}
\renewcommand{\labelenumii}{\alph{enumi}$_\arabic{enumii}$)}
\thispagestyle{empty}
\begin{center}
{\large {\bf M\'etodos Num\'ericos}}\\


\vspace*{1.5cm} 
\rule{0em}{2.em}{\bf Gu\'{i}a 1} \\ Marzo de 2024
\end{center}
\vspace*{.5cm} 


\noindent
{\bf Nota:} La informaci�n, en la computadora que use, es guardada
de forma permanente, en lo que se suele llamar disco r�gido.
Es su responsabilidad mantener un sistema ordenado de trabajo,
donde la informaci�n est� guardada en un �rbol natural
de organizaci�n.
Dependiendo del sistema operativo, cuando se inicia una ventana
de trabajo o terminal, se puede estar en la ubicaci�n inicial,
que gen�ricamente llamaremos \emph{home}, 
que por ejemplo podr�a ser:


{\tt /home/jaimito}

si se est� en linux (ubuntu).
En el lugar donde en el ejemplo aparece {\tt jaimito}
deber�a aparecer su nombre de usuario.
De aqu� en adelante asumimos que trabaja con un sistema operativo linux,
como por ejemplo \emph{ubuntu}.



%\bigskip

%\newpage

\noindent {\bf Ejercicios de linux}


%1
\prob{} En una terminal, estando en su directorio inicial de trabajo (home),
cree un directorio nuevo llamado {\tt metodos\_numericos};
a prop�sito sin acento. 
Cambie de directorio desde el inicial al directorio
{\tt metodos\_numericos}.
En ese directorio cree otro directorio llamado {\tt guias};
a prop�sito sin acento. 
Descargue desde el
navegador el archivo de esta gu\'ia y mueva el mismo al nuevo directorio.
Cambie el nombre del archivo a {\em guia1.pdf}\, y copie el archivo a otro 
con nombre {\em Copia.pdf}.  Luego, liste el contenido del
directorio y verifique los nombres y el tama\~no y fecha de los mismos.
Averig\"ue qu\'e espacio ocupa el directorio usando el comando {\tt du -h}.
Finalmente, borre el archivo {\em Copia.pdf}.
Para ver en qu� lugar del �rbol de directorio est�, ejecute el comando
{\tt pwd}.


%2
\prob{} 
\begin{itemize}

\item Aseg\'urese de que se encuentra en el directorio {\tt metodos\_numericos}  (pwd).
 
\item Cree 2 carpetas: ejercicios1, ejercicios2 (utilizando ``mkdir'').
 
\item Acceda a ejercicios1 (utilizando ``cd'').
 
\item All\'i, cree un archivo de texto llamado "file1.dat" con la frase "Hola mundo" (utilizando ``kate'').
Aseg\'urese de guardarlo.
 
\item Sin moverse del  directorio ejercicios1, haga una copia del archivo "file1.dat" dentro de ejercicios2, 
y ll\'amela "file1\_copia.dat" (utilizando ``cp'').
 
\item Acceda al  directorio ejercicios2 (con ``cd'') y verifique la existencia del archivo "file1\_copia.dat" 
en ese directorio (con ``ls'').
 
\item Muestre el  contenido del archivo usando el  comando "cat".
 
 
\item Cambie el nombre del archivo "file1\_copia.dat" a "file1\_movido.dat", utilizando ``mv''.
 
 
\item Ub\'iquese en el directorio {\tt metodos\_numericos}, haciendo ``cd ..''. Haga ``ls'' 
(para ver el contenido) y ``pwd''. 

\item Borre el directorio ``ejercicios2'' con todo su contenido, haciendo ``rm -Rf ejercicios2''.
 
\end{itemize}



%\bigskip

%\newpage

\noindent {\bf Ejercicios de gnuplot}

%3
\prob{} Utilice el programa {\em gnuplot} para graficar la funci\'on
$f(x) = \sin (2\pi x)$, para $x  \in [0:5]$. Ajuste la cantidad de puntos
que usa gnuplot para graficar ({\em samples}) a 500 para mejorar la calidad
del gr\'afico. D\'e nombre a los ejes, t\'itulo al gr\'afico y defina la
leyenda para que aparezca $f(x)$.

%4
\prob{} Abra una terminal y realice las siguientes acciones:
\begin{enumerate}
\item 
Vaya al directorio:

{\tt /home/jaimito/metodos\_numericos/}

Donde, en lugar de ``jaimito'', debe ser el nombre de su usuario. All� cree el directorio

{\tt guia1}


\item En ese nuevo directorio ejecute el editor {\tt kate} liberando
la l�nea de comandos (utilice el s�mbolo $\&$).

\item En el editor de texto ingrese los siguientes datos en tres columnas,
separados por dos espacios:
\begin{center}
\begin{tabular}{ccc}
0.0  & 1.0  & 1.2\\
1.0  & 1.5  & 1.6  \\
2.0  & 2.0  & 2.0 \\
3.0  & 2.5  & 2.4 \\  
\end{tabular}
\end{center}
\noindent
Grabe esta tabla en el archivo {\em datos.dat}. 
%Aseg\'urese que est\'e en el directorio {\em MetNum}. 
No cierre el editor. 

\item Ahora utilice {\tt gnuplot} para graficar la 2da y la 3era columna versus
la primera. Grafique usando s\'olo puntos. En el mismo gr\'afico incluya las funciones 
$y = 0.5 x + 1$ y $y=0.4 x + 1.2$. 
Nombre al eje horizontal {\em x} y al vertical {\em y}. Habilite la grilla. 
Titule el gr\'afico {\em Funciones lineales}. Cree archivos
de salida en formato png, pdf y postscript color y blanco y negro (n\'ombrelos 
p4.png, p4.pdf, p4\_color.ps y p4\_byn.ps respectivamente). 
Abra otra terminal y con la aplicaci\'on {\tt  okular} visualice
los archivos de salida, a medida que los va generando). Luego, experimente graficar 
cambiando el tama\~no de los puntos y el grosor de las l\'ineas, etc. 


\item Liste los archivos del directorio mostrando su tama\~no y fecha. 
Averig�e qu\'e espacio ocupan los archivos.
Finalmente, cierre la ventana del editor. 
Borre los archivos postscript y cierre la terminal.

\end{enumerate}



%\bigskip

%\newpage

\noindent {\bf Ejercicios de Fortran}

%5
\prob{} 
Evaluar con un programa {\sc Fortran} las siguientes operaciones
matem�ticas. Realice todas las operaciones en forma secuencial en un mismo programa.
Imprima por pantalla cada operaci\'on y su resultado. 

\begin{enumerate}
\item {\tt A =  5 / 2 + 20/ 6} 
\item {\tt B =  4 * 6 / 2 - 15 / 2}
\item {\tt C = 5 * 15 / 2 / (4 - 2)} 
\item {\tt D = 1 + 1/4}
\item {\tt E = 1. + 1/4}
\item {\tt F = 1 + 1./4}
\item {\tt G = 1. + 1./4.}
\end{enumerate}

%6
\prob{} Escriba un programa {\sc Fortran}
que pida dos n\'umeros reales e imprima en la
pantalla el mayor de ellos. El programa debe indicar si los n\'umeros son iguales.

%7
\prob{} Escriba un programa {\sc Fortran}
que pida un n\'umero entero y determine si es
m\'ultiplo de 2 y de 5.



%8
\prob{} Escriba una programa que ingrese los coeficientes $A$, $B$ y $C$ de un
polinomio real de segundo grado ($A x^2 + B x + C$), calcule e imprima en
pantalla las dos ra\'ices del polinomio en formato complejo $a + i b$, sin
utilizar algebra compleja.




%9
\prob{} Implementar un programa Fortran para evaluar la suma (en precisi\'on
simple)
\[
\sum_{n=1}^{10\,000\,000} \frac{1}{n}
\]
\noindent
primero, en el orden usual, y luego, en el orden opuesto. Explique las diferencias
obtenidas  e indique cu\'al es m\'as preciso y su justificaci\'on.



%10
\prob{} Dise\~nar una funci\'on que calcule la potencia en\'esima de un n\'umero, es decir que devuelva $X^n$ para $X$ real y $n$ entero. Realice un programa que utilice la funci\'on e imprima en pantalla las primeras 5 potencias naturales de un n\'umero ingresado. 

Usar m�dulos, es decir, utilizar un archivo ``parametros.f90'' para definir la precisi�n, otro archivo ``funciones.f90'', para la funci�n y otro archivo para el programa principal. Emplear un script de compilaci�n (``compilar.sh''). El mismo contendr� los comandos necesarios para compliar todo:

gfortran parametros.f90 funciones.f90 prob10.f90 -o prob10.x

y deber� ser un archivo ejecutable. Para esto, en la consola, escribir: chmod u+x compilar.sh




%11
\prob{} Escriba un programa, que utilizando una subrutina, multiplique un
vector de $N$ elementos, por una matriz de $N\times N$. El programa debe
preguntar el valor de $N$ y luego definir los arreglos, y
darle valores iniciales tal que, la matriz sea triangular superior, con todos
sus elementos igual a 1, excepto los de la diagonal que toman valor 3.
El vector tendr\'a todos sus elementos pares igual a 2, y los impares igual a 3. 
Utilizar m�dulos.


%12
\prob{}
Se determinaron 9 mediciones de concentraciones de f�sforo en un campo y los 
datos obtenidos fueron:

479.80 ; 499.10 ; 510.20 ; 508.60; 503.30 ; 501.20; 507.30 ; 499.90 y 468.60.

\begin{enumerate}

\item Escriba estos valores en un archivo, en forma de columna y llame a este archivo 
\textit{datos1.dat}.

\item Escriba un programa en fortran que pregunte el n\'umero de datos ($9$ en este caso), 
a ser ingresado por teclado, que lea del archivo \textit{datos1.dat} todos los valores
guard\'andolos en un vector y que luego calcule y muestre por pantalla el promedio 
y la desviaci\'on est\'andard de los datos. Adem\'as de mostrar el resultado por pantalla, 
el programa debe escribirlo en un archivo llamado \textit{resultados1.dat}. 
El programa debe llamarse \textit{prob12.f90}

\item Utilice el programa para calcular el promedio y  desviaci\'on est\'andard de los datos
mencionados.


\end{enumerate}

Nota: 

El promedio se calcula de la siguiente forma:

\begin{equation}
 \overline{x} = \frac{1}{n} \sum_{i=1}^{n} x_i
\end{equation}

\noindent siendo $n$ el n\'umero de datos.


La  desviaci\'on est\'andard se calcula de la siguiente forma:

\begin{equation}
 \sigma = \sqrt{ \frac{1}{n-1} \sum_{i=1}^{n} ( x_i - \overline{x})^2 }
\end{equation}

\noindent siendo $n$ el n\'umero de datos.

Alternativamente se puede utilizar la siguiente f\'ormula:


\begin{equation}
\sigma = \sqrt{ \frac{1}{(n-1)} ( \sum_{i=1}^{n}  x_i^2 -  n \overline{x}^2 )}
\end{equation}




%13
\prob{} Escriba un programa que permita convertir n\'umeros naturales con base 
$10$ a la base $b <= 16$. El programa
debe pedir como entrada $b$ y el n\'umero natural a convertir (que debe estar en
base 10). 
Utilice {\tt SELECT CASE} para asignar los d\'igitos (del resultado en base $b$). 
Utilice {\tt CASE DEFAULT} (para asignar los d\'igitos del 0 al 9).




%14
\prob{} Efect\'ue con un programa en Fortran en simple precisi\'on los siguientes
c\'alculos, matem\'aticamente equivalentes,

\begin{enumerate}
\item $1\,000\,000 \times 0.1$

\item $\sum_{n=1}^{1\,000\,000} 0.1$

\item $\sum_{m=1}^{1\,000} \left( \sum_{n=1}^{1\,000} 0.1\right)$


\end{enumerate}
\noindent

{ Explique las diferencias obtenidas entre resultados finales de a), b) y c) y muestre que el error relativo en b) es del orden
del 1\%, pero es mucho menor en c). Resalte la conclusi\'on de este ejercicio. }


\begin{enumerate}
\item[{\bf d)}] En los puntos b) y c), vaya guardando con un write (cada 1000 iteraciones en el caso b, 
y en todas las iteraciones de la suma externa en c), los errores parciales en un archivo de datos.  
La primera columna con los valores de los \'indices (multiplicados por 1000 para el caso c) de las sumas 
y segunda columna con el error parcial, i.e., la diferencia entre el valor de la respectiva 
suma parcial y el valor exacto ($i*0.1$).  
Grafique usando {\em gnuplot}, los errores parciales de b) y c) superpuestas, en funci\'on del 
\'indice correspondiente, en un archivo``prob14.pdf" en color, con t\'itulo de gr\'afica``Problema 14, 
Gu\'ia 1", con nombres adecuados en los ejes y con leyendas adecuadas. 
Los datos del b), graficarla con puntos de tama\~no 0.75 (el default es 1), y los de c) con l\'inea
de grosor 1.6 (el default es tambi\'en 1). 
Luego de ver c\'omo le queda la gr\'afica en escala lineal, cambie el eje x, 
a escala logar\'itmica. Seg\'un su criterio, en qu\'e escala se 
aprecia mejor el resultado del problema para su an\'alisis? Puede al eje $y$ darle escala 
logar\'itmica? Pruebe qu\'e pasa si lo hace.
Analice la conveniencia en general del uso de escalas logar\'itmicas.

Finalmente deje el archivo grabado en escala lineal en $y$, y logar\'itmica en $x$, y diga el 
tama\~no del mismo tanto en bytes como Kbytes. 


\end{enumerate}
 

%15
\prob{}

Calcular sin(x), seg\'un su desarrollo en series de Taylor, teniendo en cuenta que:


\begin{equation}
 sen(x) = \sum_{n=0}^{\infty} \frac{(-1)^n}{(2n+1)!} x^{2n+1} 
\end{equation}

Aproximar con los $9$ primeros t\'erminos (hasta $n=8$).
Es decir, calcular:


\begin{equation}
 sen(x) = \sum_{n=0}^{8} \frac{(-1)^n}{(2n+1)!} x^{2n+1} 
\end{equation}

Para ello:

\begin{enumerate}

\item Resuelva el problema para $x=0.1$.

\item Realice los c\'alculos en simple precisi\'on.

\item Guarde en un vector de n\'umeros enteros, las factoriales hasta $2 n_{max} + 1$, 
siendo $n_{max}=8$. (Ayuda1: Utilice un ``do'' para ello. Ayuda2: al realizar operaciones, 
preste atenci\'on a si est\'a trabajando con enteros o con reales).

\item Calcule la suma desde $0$ hasta $n_{max}$ de los t\'erminos en orden directo.

\item Calcule la suma  de los t\'erminos en orden inverso (es decir, desde $n_{max}$ hasta $0$).

\item Muestre por pantalla el resultado de ambas sumas y el verdadero valor de $sin(x)$.

\item Para el caso de la suma directa, vaya guardando en un archivo (llamado $seno.dat$) 
los sucesivos pasos de la suma y el error cometido (la diferencia con el verdadero 
valor de $sin(x)$).

\item Analice cu\'al de las dos formas de sumar se aproxima m\'as 
al resultado exacto y explique por qu\'e.

\end{enumerate}



\bigskip

\bigskip

%\newpage

\noindent {\bf Ejercicios Complementarios}


%16
\prob{} Escriba un programa para calcular la posici\'on y la 
velocidad en funci\'on del tiempo, para un problema de tiro oblicuo.
Debe preguntar el \'angulo (en grados) y la velocidad 
inicial (en m/seg.), asumiendo que el proyectil parte del origen. 
Elija el incremento temporal ($\Delta t$) de manera que la gr\'afica
tenga 600 puntos y abarque el intervalo entre el disparo y el instante en
que el proyectil vuelve a tener altura 0. Utilice {\em funciones} para calcular
la posici\'on y la velocidad. Escriba la salida del programa en un archivo
de texto, con 5 columnas ($t$, $x(t)$, $y(t)$,$v_x(t)$,$v_y(t)$). La primera
l\'inea debe comenzar con \#, e incluir la descripci\'on de los datos
de cada columna. Los datos en la tabla deben tener 6 cifras significativas y
estar escritos en notaci\'on exponencial. Grafique $x(t)$, $y(t)$ y $v_y(t)$ 
en funci\'on de $t$, y la trayectoria del proyectil, utilizando {\em gnuplot}.






%17 
\prob{} Escribir un programa que pida una contrase\~na de tres d\'igitos y permita leer tres intentos.  Si el usuario da la contrase\~na correcta responde
``Correcto'' y queda inactivo, con este mensaje.  En caso contrario el programa
escribe ``Lo siento, contrase\~na equivocada'' y se cierra de inmediato.

%18 
\prob{} Escribir un programa que, dado un a\~no y el nombre de un mes, 
saque por pantalla el n\'umero de d\'ias del mes (tenga en cuenta que 
algunos a\~nos son bisiestos). 



%19
\prob{} Escriba un programa para calcular un valor aproximado de $\pi$
utilizando 
\begin{enumerate}
\item la f\'ormula recurrente de Arqu\'imedes, que  acota $\pi$ entre
$P_n$ y $p_n$, con $p_n < \pi < P_n$. Siendo $p_n$ y $P_n$ 
los per\'imetros de los pol\'igonos regulares de $n$ lados inscriptos y 
circunscriptos, respectivamente, en la circunsferencia de radio $1/2$. 
La f\'ormula de recurrencia que encontr\'o Arqu\'imides es la siguiente:
\begin{align*}
P_{2n} = \frac{2 p_n P_n}{p_n + P_n} \\
p_{2n} = \sqrt{P_{2n} p_n}
\end{align*}
Usando los valores $P_6=2\sqrt{3}$ y $p_6 = 3$, correspondientes al hex\'agono, 
escriba un programa que realice 20 iteraciones, con $n = 6\times 2^k$, y 
$k=1, \cdots, 20$, y escriba los resultados en pantalla.

\item la productoria de Wallis
\[
\frac{\pi}{2} = \prod_{n=1}^\infty \frac{(2n)^2}{(2n)^2-1} =
\frac43\frac{16}{15}\frac{36}{35}\frac{64}{63}\dots 
%\frac21\frac23\frac43\frac45\frac65\frac67\dots 
\]
Calcule el valor de $\pi$ truncando la productoria a $10^6$ factores.

\end{enumerate}



%20
\prob{} Escriba el siguiente programa y explique por qu\'e los valores
obtenidos no son iguales: \\ 

{\fontfamily{phv}\selectfont
\hspace*{4cm}   program test  \\
\hspace*{4cm}   implicit none  \\
\hspace*{4cm}   integer, parameter   \hspace*{0.1cm}  :: pr = kind(1.0)\\
\hspace*{4cm}   real(pr)         \hspace*{1.8cm}       :: sum0,sum1 \\
\hspace*{4cm}   integer          \hspace*{1.9cm}     :: i  \\
\hspace*{4cm}   ! \\
\hspace*{4cm}   sum0 = 0.\_pr ; sum1 = 1.\_pr \\
\hspace*{4cm}   do i=1,10000 \\
\hspace*{4cm}   \hspace{.3cm} sum0 = sum0 + 1.e-8\_pr \\
\hspace*{4cm}   \hspace{.3cm} sum1 = sum1 + 1.e-8\_pr \\ 
\hspace*{4cm}   end do \\
\hspace*{4cm}   sum0 = sum0 + 1.\_pr  \\
\hspace*{4cm}   write(*,*) sum0, sum1  \\
\hspace*{4cm}   end program test }


%21
\prob{} Dado el siguiente programa y su correspondiente subrutina:

\vspace{.3cm}

\begin{minipage}{0.5\textwidth}
{\fontfamily{phv}\selectfont
\hspace*{1cm}   program test  \\
\hspace*{1cm}   implicit none  \\
\hspace*{1cm}   real(kind(1.)) :: x, y,s  \\
\hspace*{1cm}   integer        :: i,j  \\
\hspace*{1cm}   i = 6; j = 3; x = 4.; y = 8.   \\
\hspace*{1cm}   call sum(i,x,s)  \\
\hspace*{1cm}   write(*,'(F4.1," + ",F4.1," = ",F5.1)') real(i),x,s   \\
\hspace*{1cm}   end program test}
\end{minipage}\hfill
\begin{minipage}{0.5\textwidth}
{\fontfamily{phv}\selectfont
\hspace*{2cm}   \phantom{a}  \\
\hspace*{2cm}   subroutine sum(z,w,ss)  \\
\hspace*{2cm}   implicit none  \\
\hspace*{2cm}   real(kind(1.)), INTENT(IN) \hspace{.3cm}  :: z, w  \\
\hspace*{2cm}   real(kind(1.)), INTENT(OUT) :: ss  \\
\hspace*{2cm}   ss = z + w  \\
\hspace*{2cm}   end subroutine sum   } \\
\end{minipage}

\vspace{.3cm}
\noindent
Escriba, compile y ejecute el programa de las siguiente manera:
\begin{enumerate}
\item con el programa tal como est\'a, incluyendo la subrutina en el mismo archivo o en 
archivo separado.

\item modificando el programa con la cl\'ausula {\tt CONTAINS}, para que {\em contenga} a
la subrutina.

\item creando un modulo que contenga la subrutina, y usando este en el programa.

\end{enumerate}

\noindent
Verifique que en el primer caso compila sin errores y produce {\em resultados incorrectos},
y en los otros casos no.




%22
\prob{} Interprete el resultado de los siguientes programas en virtud de la representaci\'on
de punto flotante de los n\'umeros reales.\\
\begin{itemize}
\item 
{\fontfamily{phv}\selectfont
\hspace*{4cm}  program test\textunderscore igualdad  \\
\hspace*{4cm}  implicit none  \\
\hspace*{4cm}  integer, parameter   \hspace*{0.5cm}  :: pr=kind(1.d0)  \hspace*{0.5 cm} ! pr puede ser simple o doble   \\ 

\hspace*{4.2cm}  if (19.08\textunderscore pr + 2.01\textunderscore pr == 21.09\textunderscore pr) then  \\
\hspace*{4.5cm}     write(*,*) '19.08 + 2.01 = 21.09 '  \\
\hspace*{4.4cm}  else  \\
\hspace*{4.6cm}     write(*,*) '19.08 + 2.01 /= 21.09 '  \\
\hspace*{4.2cm}  endif  \\
\hspace*{4cm}  end program test\textunderscore igualdad }  \\

\vspace{1cm}


\item 
{\fontfamily{phv}\selectfont
\hspace*{4cm}  program test\textunderscore igualdad2  \\
\hspace*{4cm}  implicit none  \\
\hspace*{4cm}  integer, parameter   \hspace*{0.5cm}  :: pr=kind(1.d0)  \hspace*{0.5 cm} ! pr puede ser simple o doble   \\  
\hspace*{4cm}  real(pr)             \hspace*{0.75cm} :: a  \\ 

\hspace*{4.2cm}  a = 2.05\textunderscore pr  \\
\hspace*{4.2cm}  if (a*100.\textunderscore pr  == 205.\textunderscore pr) then  \\
\hspace*{4.5cm}     write(*,*) '2.05*100 = 205 '  \\
\hspace*{4.4cm}  else  \\
\hspace*{4.6cm}     write(*,*) '2.05*100 /= 205 '  \\
\hspace*{4.2cm}  endif  \\
\hspace*{4cm}  end program test\textunderscore igualdad2 }  \\

\end{itemize}
%Piense el mensaje de este ejercicio, el cual debe tener presente en toda la materia.\\




%23
\prob{} Escriba un programa, que utilizando una subrutina, multiplique un
vector de $N$ elementos, por una matriz de $N\times N$. El programa debe
preguntar el valor de $N$ y luego definir los arreglos, y
darle valores iniciales tal que, la matriz sea triangular superior, con todos
sus elementos igual a 1, excepto los de la diagonal que toman valor 3.
El vector tendr\'a todos sus elementos pares igual a 2, y los impares igual a 3. 
No utilice {\tt DO} para las inicializar el vector, ni {\tt DO} anidados para
inicializar la matriz. 


%24
\prob{} La f\'ormula cuadr\'atica nos dice que las ra\'ices de $ax^2 + bx + c = 0$
son
$$
x_1 = \frac{-b + \sqrt{b^2 -4ac}}{2a}, \qquad \qquad x_2 =  \frac{-b - \sqrt{b^2 -4ac}}{2a} .
$$
\noindent
Si $b^2\gg 4ac$, entonces, cuando $b>0$ el c\'alculo de $x_1$ involucra en el
numerador la sustracci\'on de dos n\'umeros casi iguales, mientras que si $b<0$,
esta situaci\'on ocurre para el c\'alculo de $x_2$. ``Racionalizando el numerador''
se obtienen las siguientes f\'ormulas alternativas que no sufren este problema:
$$
x_1 = \frac{-2c}{b + \sqrt{b^2 -4ac}}, \qquad \qquad x_2 =  \frac{2c}{-b + \sqrt{b^2 -4ac}} ,
$$
\noindent
siendo la primera adecuada cuando $b>0$, y la segunda cuando $b<0$.
Escriba un programa en precisi\'on simple que utilice la f\'ormula usual y la 
``racionalizada'' para calcular las ra\'ices de
$$
x^2 + 6210 x + 1 = 0.
$$
\noindent 
Interprete los resultados.

\bigskip

 
\end{document}

